\documentclass[12pt, a4paper]{academics}

\academicsCoverLabel{解题报告}
\academicsCourse{算法设计与分析}{Design and Analysis of Algorithms}
\academicsReport{序列动态规划}{2026 年 6 月 4 日}
\academicsArthur{华博文}{19240212}

\begin{document}

\makeacademicscover

\section{引言}

动态规划是算法设计与分析中最为核心的方法论之一。其本质在于将原问题拆解为相互关联的子问题，利用子问题的最优解递推地构造原问题的最优解。这一思想看似朴素，但在实际运用中，状态如何定义、转移如何设计、复杂度瓶颈如何突破，每一步都考验着解题者对问题结构的洞察力。

本报告以最长上升子序列（Longest Increasing Subsequence，以下简称 LIS）为主线，依次展开四道关联紧密的题目。从最基础的 $\mathrm{O}(n^2)$ 动态规划入门，到借助单调性将复杂度优化至 $\mathrm{O}(n \log n)$，再到将二维偏序约束通过排序技巧化归为一维 LIS，最终将其与最长公共子序列相融合，导出经典的 LCIS 问题及其 $\mathrm{O}(nm)$ 解法。四道题目层层递进，每道题在前一题的基础上引入一个新的维度，既保持了算法思想的连贯性，又迫使解题者在每一阶段重新审视问题结构、寻找更深刻的优化方向。

报告中的所有代码均使用 C++23 标准编写，在 Apple clang 21.0.0 编译器下通过测试。

\section{题一：最长上升子序列}

\subsection{问题重述}

给定一个长度为 $n$ 的整数序列 $a_1, a_2, \dots, a_n$，需要求出其中最长严格上升子序列的长度。所谓子序列，是指从原序列中删除若干元素（也可以不删除任何元素）后，保持原有顺序得到的序列；严格上升则要求子序列中每一个元素都严格大于前一个元素。本题的数据范围为 $1 \le n \le 1000$，序列中元素的绝对值不超过 $10^9$。

\subsection{状态设计与转移方程}

子序列类问题天然适合用动态规划求解，其原因在于问题具备最优子结构性质：以某个位置结尾的最长上升子序列，必然由一个更短的、以某个更小元素结尾的上升子序列拼接上当前元素而构成。这提示我们以「结尾位置」作为 DP 状态的划分依据。

具体而言，设 $\mathrm{dp}[i]$ 表示以 $a_i$ 为结尾的最长严格上升子序列的长度。对于每个位置 $i$，我们考察所有排在它之前的位置 $j$（$j < i$），如果 $a_j < a_i$，则 $a_i$ 可以接在以 $a_j$ 结尾的上升子序列之后，形成一个长度增加 $1$ 的新子序列。由此得到状态转移方程：

\[
\mathrm{dp}[i] = \max_{1 \le j < i, a_j < a_i}\{\mathrm{dp}[j] + 1\}
\]

每个元素自身至少构成长度为 $1$ 的子序列，因此 $\mathrm{dp}[i]$ 的初始值均为 $1$。最终答案为所有 $\mathrm{dp}[i]$ 中的最大值，即 $\displaystyle \max_{1 \le i \le n} \mathrm{dp}[i]$。

这一解法的时间复杂度为 $\mathrm{O}(n^2)$，空间复杂度为 $\mathrm{O}(n)$。对于 $n \le 1000$ 的数据规模而言，至多 $10^6$ 次运算在毫秒级别内即可完成，完全满足题目要求。然而，当数据范围扩展至 $10^5$ 时，$10^{10}$ 次运算将远远超出时限，这使得我们必须重新审视问题的结构，寻找更为高效的策略——这便是下一道题目的出发点。

\subsection{代码实现}

\inputminted{c++}{../src/p1_lis_dp.cpp}

\section{题二：最长上升子序列的优化}

\subsection{问题重述}

本题的问题描述与题一完全一致，唯一的变化在于数据范围扩大至 $1 \le n \le 10^5$。$\mathrm{O}(n^2)$ 的朴素动态规划在此规模下已不可行，需要寻求时间复杂度更低的解法。

\subsection{贪心策略与单调性分析}

回顾题一的 DP 过程，我们为每个位置 $i$ 维护了以 $a_i$ 结尾的最优解。然而，随着序列长度的增长，记录每一个位置的 $\mathrm{dp}$ 值不仅带来了平方级别的状态转移开销，更关键的是，我们实际上并不关心所有可能的结尾——在长度相同的前提下，结尾元素越小，后续元素能够接在其后的可能性就越大。这一观察指向了一个贪心思路：对于每种长度，只记录具有最小结尾元素的那个子序列。

形式化地，我们维护一个数组 $\mathit{tails}$，其中 $\mathit{tails}[k]$ 表示在当前已处理的元素中，所有长度为 $k+1$ 的上升子序列里，结尾元素的最小可能值。可以证明，$\mathit{tails}$ 数组始终严格单调递增。这一性质的正确性可以通过反证法加以确认：若存在 $k < \ell$ 使得 $\mathit{tails}[k] \ge \mathit{tails}[\ell]$，则从长度为 $\ell+1$ 的对应子序列中截取前 $k+1$ 个元素，便得到一个长度为 $k+1$ 且结尾元素严格小于 $\mathit{tails}[\ell] \le \mathit{tails}[k]$ 的上升子序列，这与 $\mathit{tails}[k]$ 的最小性定义矛盾。

\subsection{算法流程}

基于 $\mathit{tails}$ 的单调性，我们可以为每个新元素 $a_i$ 在 $\mathit{tails}$ 中执行一次二分查找，定位第一个大于或等于 $a_i$ 的位置 $\mathit{pos}$。若 $\mathit{pos}$ 存在，则意味着 $a_i$ 能够以更小的代价取代原有的结尾元素，我们将 $\mathit{tails}[\mathit{pos}]$ 更新为 $a_i$；若 $\mathit{pos}$ 不存在，即 $a_i$ 严格大于 $\mathit{tails}$ 中的所有元素，则将 $a_i$ 追加到 $\mathit{tails}$ 末尾，意味着我们找到了一个更长的上升子序列。

这一算法在文献中被称为 Patience Sorting，得名于扑克牌游戏中的分堆策略。值得注意的是，替换操作并不会丢失任何有用的信息：由于 $\mathit{tails}$ 中各位置的元素在替换后只可能变小而不会变大，在替换之后，原本可以用更大结尾元素完成的扩展仍然可行，甚至因为结尾的变小而拥有了更多的扩展机会。而更短长度的子序列信息则由 $\mathit{tails}$ 中更靠前的位置所保留，不受后续替换的影响。

遍历结束后，$\mathit{tails}$ 的最终长度即为最长上升子序列的长度。由于每个元素执行一次 $\mathrm{O}(\log n)$ 的二分查找，总时间复杂度为 $\mathrm{O}(n \log n)$，空间复杂度为 $\mathrm{O}(n)$。对于 $n = 10^5$ 的规模，约 $1.7 \times 10^6$ 次运算可以在数十毫秒内完成。

\subsection{代码实现}

\inputminted{c++}{../src/p2_lis_fast.cpp}

\section{题三：信封嵌套问题}

\subsection{问题重述}

现有 $n$ 个信封，每个信封由宽度 $w_i$ 和高度 $h_i$ 两个属性刻画。若信封 $A$ 的宽度和高度均严格大于信封 $B$，则称 $A$ 可以套住 $B$。一个信封内部可以嵌套多个信封，如同俄罗斯套娃一般。需要求出最多能嵌套的信封数量，即最长的嵌套链长度。数据范围为 $1 \le n \le 10^5$，宽和高均不超过 $10^9$。

\subsection{从二维偏序到一维 LIS}

本题与经典 LIS 的关键区别在于约束条件从单一维度扩展为两个维度。对于一个合法的嵌套链，前一信封的宽度和高度必须同时严格小于后一信封，这构成了一个典型的二维严格偏序关系。如果能够将其中一个维度「消去」，就可以将问题化归为我们已经解决的 LIS 问题。

降维的关键在于排序策略。我们首先将所有信封按宽度升序排列。经过这一处理后，在任意一个前缀中，后续信封的宽度天然不小于之前的信封，宽度这一维度的顺序约束便得到了满足。然而，宽度相同的信封之间并不允许嵌套——如果直接对宽度相同的一组信封按高度升序排列，LIS 算法有可能错误地将同宽的信封选入嵌套链中。

解决这一矛盾的方法颇具巧思：在宽度相同的情况下，对高度进行降序排列。如此一来，所有同宽信封在排序后依次出现，且高度递减。LIS 算法要求子序列严格递增，因此在面对高度递减的序列时，至多只能选取其中的一个元素（事实上，只有第一个被处理到的具有合适高度的同宽信封有可能被纳入 LIS），从而天然规避了同宽嵌套的非法情况。

排序完成后，宽度维度的约束已被排序所吸收，我们只需在提取出的高度序列上执行题二的 $\mathrm{O}(n \log n)$ 严格 LIS 算法即可。排序本身耗时 $\mathrm{O}(n \log n)$，LIS 同样耗时 $\mathrm{O}(n \log n)$，总体时间复杂度仍为 $\mathrm{O}(n \log n)$，空间复杂度为 $\mathrm{O}(n)$。这一「排序降维、LIS 求解」的策略在二维及更高维度的偏序问题中均有广泛的应用，典型的例子还包括俄罗斯套娃信封问题在三维空间中的推广——此时则需要借助 CDQ 分治或树状数组等更复杂的工具。

\subsection{代码实现}

\inputminted{c++}{../src/p3_envelopes.cpp}

\section{题四：最长公共上升子序列}

\subsection{问题重述}

给定两个整数序列 $A$（长度为 $n$）和 $B$（长度为 $m$），需要求出一个序列，使得它既是 $A$ 的子序列，又是 $B$ 的子序列，且严格单调递增。在所有满足条件的序列中，求出长度的最大值。数据范围 $1 \le n, m \le 3000$。

\subsection{复合约束下的状态设计}

本题同时融合了最长公共子序列（LCS）和最长上升子序列（LIS）两个经典问题的约束，属于二者的交叉变体，在算法竞赛中通常被称为 LCIS（Longest Common Increasing Subsequence）。表面上看，我们似乎可以分别处理公共性和递增性——但这两者相互缠绕，难以解耦。若沿用 LCS 的状态定义 $\mathrm{dp}[i][j]$ 表示 $A[1..i]$ 与 $B[1..j]$ 的 LCS 长度，则无法从中获知子序列最后一个元素的值，自然也无从保证后续元素的递增性；若只沿用 LIS 的状态设计，则丢失了公共子序列的约束。

解决这一困境的思路在于，将递增性约束直接编码到状态之中——在 DP 状态中不仅记录长度，还要明确记录「以哪个元素结尾」。考虑到公共性要求选出的元素必须同时出现在 $A$ 和 $B$ 中，最自然的选择是以 $B$ 中的某个元素作为结尾的锚点。由此，我们定义 $\mathrm{dp}[i][j]$ 表示在 $A[1..i]$ 和 $B[1..j]$ 的范围内，以 $b_j$ 为最后一个元素的公共上升子序列的最大长度。

这一状态设计有两个精妙之处。其一，结尾元素固定为 $b_j$，使得我们在考虑后续扩展时可以直接比较元素大小，从而保证了递增性的可判定。其二，状态的前一维 $i$ 是可选的——$\mathrm{dp}[i][j]$ 并不强制要求 $a_i$ 被选入，如果 $a_i$ 无法与 $b_j$ 配对，则直接继承 $\mathrm{dp}[i-1][j]$ 的值，这与传统 LCS 的转移逻辑一脉相承。

\subsection{转移方程}

在遍历过程中，对每个固定的 $i$，我们依次考察 $B$ 中的每个位置 $j$。根据 $a_i$ 与 $b_j$ 的关系，分两种情况讨论。

当 $a_i \neq b_j$ 时，$b_j$ 无法与当前的 $a_i$ 配对，因此以 $b_j$ 结尾的 LCIS 只能来源于 $A$ 的前 $i-1$ 个元素，即 $\mathrm{dp}[i][j] = \mathrm{dp}[i-1][j]$。

当 $a_i = b_j$ 时，$b_j$（也即 $a_i$）可以作为 LCIS 的新结尾。它必须接在 $B$ 中某个更靠前、且值更小的元素 $b_k$ 之后（$k < j$ 且 $b_k < b_j$）。于是转移方程为：

\[
\mathrm{dp}[i][j] = \max_{\substack{1 \le k < j \\ b_k < b_j}}\{\mathrm{dp}[i-1][k]\} + 1
\]

\subsection{滚动维护最值的优化}

上述转移方程在朴素实现下需要对每个 $(i, j)$ 枚举 $k$，总复杂度为 $\mathrm{O}(nm^2)$，在 $n, m \le 3000$ 时达到 $2.7 \times 10^{10}$ 级别，完全不可接受。

仔细观察可以发现，对于固定的 $i$ 和遍历中的 $j$，内层 $\max$ 所依赖的 $k$ 的取值范围随着 $j$ 的增加而单调扩大的——每前进一步，候选集合中至多新增一个 $b_{j-1}$。更重要的是，当我们处于 $a_i = b_j$ 的情形时，条件 $b_k < b_j$ 等价于 $b_k < a_i$，而 $a_i$ 在遍历 $j$ 的过程中是固定不变的。这意味着，满足条件的 $\max \mathrm{dp}[i-1][k]$ 可以在遍历 $j$ 的过程中用一个变量 $\mathit{best}$ 动态维护：每当遇到 $b_j < a_i$，就尝试用 $\mathrm{dp}[i-1][j]$（即上一轮对应位置的 DP 值）去更新 $\mathit{best}$；当遇到 $a_i = b_j$ 时，直接使用 $\mathit{best} + 1$ 完成转移，无需内层循环。

进一步地，由于 $\mathrm{dp}[i][j]$ 仅依赖于 $\mathrm{dp}[i-1]$ 在所有 $j$ 上的取值，我们可以将 DP 数组压缩为一维，每轮遍历 $i$ 时直接原地更新，将空间复杂度从 $\mathrm{O}(nm)$ 降至 $\mathrm{O}(m)$。最终算法的时间复杂度为 $\mathrm{O}(nm)$，在 $n = m = 3000$ 时约 $9 \times 10^6$ 次运算，完全可行。

\subsection{代码实现}

\inputminted{c++}{../src/p4_lcis.cpp}

\section{总结}

本报告以最长上升子序列为核心线索，选取了四道难度递进、思想关联的题目，系统展示了动态规划从基础建模到高阶优化的完整图景。

题一的 $\mathrm{O}(n^2)$ DP 解法是所有后续讨论的起点。它以最直观的状态定义建立了 LIS 问题的动态规划模型，算法简明、正确性自明，是初学者理解「以结尾元素定义状态」这一经典范式的入门之作。

题二突破了 DP 的框架，借助 $\mathit{tails}$ 数组的单调性，将复杂度从 $\mathrm{O}(n^2)$ 压缩至 $\mathrm{O}(n \log n)$。这一优化并非 DP 技巧的简单深化，而是算法范式的根本转变——从穷举所有可能结尾的 DP 思维，转向贪心地维护每种长度的最优结尾。它提醒我们，动态规划并非总是最优解，在问题具备特殊结构性质（如此处的单调性）时，贪心策略往往能给出更高效的答案。

题三将 LIS 从一维推广至二维偏序。通过精心设计的排序策略——宽度升序、同宽时高度降序——我们成功地将宽度维度吸收进排序，将二维问题化归为熟悉的一维 LIS。这一「排序即降维」的技巧在算法竞赛中十分常见，深刻体现了「转化问题」这一元方法论的力量：与其设计全新的算法，不如将未知问题映射为已知模型。

题四更进一步，将 LIS 与 LCS 两大经典模型加以融合。复合约束迫使我们在状态设计中同时编码结尾元素信息（以 $b_j$ 为锚点），这是解决该问题的核心突破。而内层循环的滚动最值优化——在遍历过程中用一个变量动态收集并维护所需信息——是一种通用性极强的 DP 优化范式，在各类区间 DP、背包问题以及状态转移具有单调性的场景中均有广泛的应用。

回顾四道题目的递进脉络，可以看到两条贯穿始终的方法论线索。其一，状态设计是动态规划的灵魂：从题一的「以结尾定状态」到题四的「以 $b_j$ 锚定状态」，每一次问题的升级都首先体现为状态维度的增加或重新定义。其二，优化往往来源于对冗余的剔除：题二中，我们剔除了长度相同但结尾更大的子序列；题三中，排序剔除了宽度维度；题四中，滚动变量剔除了内层循环——识别并消除冗余，正是算法优化的本质所在。

\end{document}
